% Relatorio: front-tracking bifasico no HiGFlow, e a comparacao com o VOF
%
% Compila com pdflatex + bibtex + pdflatex + pdflatex, ou latexmk -pdf main
% No Overleaf: enviar a pasta inteira; o arquivo principal e' este.
%
% Os NUMEROS vivem em dados/*.dat e sao lidos por pgfplotstable, nao digitados
% no texto -- mesma convencao do relatorio da fronteira imersa.  Regenerar uma
% medida troca o .dat e o relatorio acompanha, sem a chance de a tabela e o
% texto discordarem.
\documentclass[11pt,a4paper]{article}

\usepackage[T1]{fontenc}
\usepackage[utf8]{inputenc}
\usepackage[brazil]{babel}
\usepackage{lmodern}
\usepackage{microtype}
\usepackage{amsmath,amssymb}
\usepackage{booktabs}
\usepackage{siunitx}
\usepackage{graphicx}
\usepackage{pgfplots}
\usepackage{pgfplotstable}
\usepackage{tikz}
\usepackage{xcolor}
\usepackage[hidelinks]{hyperref}
\usepackage[margin=2.7cm]{geometry}

\pgfplotsset{compat=1.18}
\sisetup{output-decimal-marker={,}, group-separator={.}}

\newcommand{\dat}[1]{dados/#1}

\title{Front-tracking bif\'asico no HiGFlow:\\
       implementa\c{c}\~ao, verifica\c{c}\~ao e compara\c{c}\~ao com o VOF}
\author{Relat\'orio t\'ecnico}
\date{Setembro de 2026}

\begin{document}
\maketitle

\begin{abstract}
Este relat\'orio documenta a implementa\c{c}\~ao de um m\'etodo de
\emph{front-tracking} para escoamento bif\'asico no HiGFlow, sua verifica\c{c}\~ao
por or\'aculos exatos, e uma compara\c{c}\~ao direta com o m\'etodo VOF j\'a
existente no mesmo sistema. A compara\c{c}\~ao foi montada de modo que os dois
m\'etodos resolvam \emph{literalmente o mesmo problema} --- mesma malha, mesma
gota, mesmas propriedades, mesmo solver e as mesmas sondas --- de sorte que a
diferen\c{c}a medida seja do m\'etodo e n\~ao do arranjo. Os resultados confirmam
duas expectativas da literatura e contrariam uma terceira: o front-tracking
obt\'em o salto de press\~ao de Laplace com erro seis vezes menor, o VOF conserva
massa em precis\~ao de m\'aquina contra uma deriva de \num{1.3e-5}, e --- ao
contr\'ario do que se esperava --- o VOF apresenta correntes par\'asitas cerca de
dez vezes menores.

Essa terceira discrep\^ancia revelou-se defeito de implementa\c{c}\~ao, e n\~ao
propriedade do m\'etodo. Levantou-se a hip\'otese de que a causa n\~ao estava na
curvatura --- que \'e exatamente onde o front-tracking \'e melhor --- mas no
\emph{balan\c{c}o discreto} entre a for\c{c}a espalhada e o gradiente de press\~ao.
A hip\'otese foi testada remontando a for\c{c}a como $\sigma\kappa\nabla H$ com
o mesmo operador discreto que a proje\c{c}\~ao inverte, e confirmou-se nos dois
sentidos que previa: as correntes par\'asitas ca\'iram mais de quatro ordens de
grandeza \emph{e} o salto de press\~ao melhorou, em vez de piorar.
\end{abstract}

\tableofcontents

\section{Introdu\c{c}\~ao}

\subsection{Contexto: a linhagem do c\'odigo}

O HiGFlow resolve as equa\c{c}\~oes de Navier--Stokes incompress\'iveis por
diferen\c{c}as finitas sobre malhas em \'arvore n\~ao graduadas, com
interpola\c{c}\~ao sem malha nas interfaces de refinamento~\cite{sousa2019}. Ele
herda uma tradi\c{c}\~ao longa de simula\c{c}\~ao de escoamentos com superf\'icie
livre no grupo: o sistema Freeflow~\cite{castelo2000}, constru\'ido sobre o
m\'etodo GENSMAC~\cite{tome1994} e sua extens\~ao tridimensional, e a
generaliza\c{c}\~ao para escoamentos multif\'asicos por
\emph{front-tracking}~\cite{sousa2004}.

Essa \'ultima refer\^encia \'e a mais pr\'oxima deste trabalho, e vale dizer
explicitamente em que: \cite{sousa2004} implementa front-tracking sobre o
GENSMAC-3D, com a frente representada por uma malha de superf\'icie
lagrangeana acoplada a uma malha euleriana de diferen\c{c}as finitas. O que este
relat\'orio documenta \'e a constru\c{c}\~ao do an\'alogo dentro do HiGFlow,
come\c{c}ando em 2D, sobre a maquinaria de fronteira imersa que j\'a existia no
sistema.

\subsection{Front-tracking \'e alternativa ao VOF, n\~ao complemento}

Os dois m\'etodos resolvem o mesmo problema f\'isico por representa\c{c}\~oes
opostas da interface:

\begin{itemize}
  \item \textbf{VOF} (\emph{volume of fluid}) \emph{captura} a interface por um
    campo euleriano de fra\c{c}\~ao volum\'etrica advectado com o escoamento. A
    interface n\~ao existe como objeto; \'e reconstru\'ida da fra\c{c}\~ao quando
    preciso. A tens\~ao superficial entra tipicamente pela formula\c{c}\~ao de
    for\c{c}a de superf\'icie cont\'inua (CSF)~\cite{brackbill1992}.
  \item \textbf{Front-tracking} \emph{rastreia} a interface por marcadores
    lagrangeanos conectados que se movem com o fluido~\cite{tryggvason2001}. A
    interface \'e um objeto geom\'etrico expl\'icito, com conectividade, e a
    curvatura sai dessa geometria.
\end{itemize}

A decis\~ao de projeto, tomada no in\'icio e mantida, \'e que os dois ficam
\textbf{separados no c\'odigo}: n\~ao se mistura a reconstru\c{c}\~ao geom\'etrica
de um com a advec\c{c}\~ao de fra\c{c}\~ao do outro. \'E o que torna poss\'ivel a
compara\c{c}\~ao da Se\c{c}\~ao~\ref{sec:comparacao}: h\'a dois m\'etodos
independentes no mesmo sistema, sobre o mesmo solver.

\subsection{Escopo: o que est\'a feito}

O projeto est\'a organizado em fases, cada uma com um or\'aculo exato antes de a
pr\'oxima introduzir f\'isica nova. O estado no momento deste relat\'orio:

\begin{center}
\begin{adjustbox}{max width=\linewidth}
\begin{tabular}{llp{7.2cm}}
\toprule
Fase & Estado & Or\'aculo \\
\midrule
B1 & feita & advec\c{c}\~ao cinem\'atica reversivel de Rider--Kothe: a frente
             volta \`a forma inicial em $t=T$ \\
B2 & feita & gota est\'atica: o salto de press\~ao bate com a lei de Laplace
             $\Delta p = \sigma/R$, e as correntes par\'asitas ficam abaixo de um
             teto medido \\
B3 & aberta & oscila\c{c}\~ao de gota na frequ\^encia de Lamb \\
B4 & aberta & bolha subindo, \emph{benchmark} publicado \\
B5 & aberta & 3D: superf\'icie triangulada \\
\bottomrule
\end{tabular}
\end{adjustbox}
\end{center}

A ordem \'e deliberada: B1 e B2 n\~ao t\^em din\^amica acoplada completa, e \'e
neles que os defeitos de m\'aquina --- advec\c{c}\~ao, cirurgia, curvatura,
espalhamento --- aparecem baratos. Tudo que se segue foi medido; onde um teste
n\~ao exercita algo, isso est\'a dito.

\section{Formula\c{c}\~ao}

\subsection{As tr\^es mudan\c{c}as em rela\c{c}\~ao \`a fronteira imersa r\'igida}

O HiGFlow j\'a possu\'ia fronteira imersa com for\c{c}amento direto: marcadores
lagrangeanos desconexos sobre um corpo r\'igido e fixo, com o n\'ucleo
regularizado de Roma~\cite{roma1999} fazendo a transfer\^encia entre malha
euleriana e marcadores. O front-tracking \'e, estruturalmente, essa mesma
maquinaria com tr\^es mudan\c{c}as.

\paragraph{1. Os marcadores advectam.} Onde o corpo r\'igido \emph{interpola} a
velocidade para calcular a for\c{c}a que imp\~oe $\mathbf{u}=0$, aqui se interpola
a velocidade para \emph{mover} o marcador:
\begin{equation}
  \mathbf{X}^{n+1} = \mathbf{X}^{n} + \Delta t\, \mathbf{u}(\mathbf{X}^{n}).
\end{equation}
\'E a mesma interpola\c{c}\~ao, com destino oposto.

\paragraph{2. A for\c{c}a \'e de tens\~ao superficial.} No lugar do multiplicador
de Lagrange do corpo r\'igido --- que vale o que for preciso para impor a
velocidade, sem lei constitutiva --- entra uma for\c{c}a \emph{constitutiva},
\begin{equation}
  \mathbf{F} = \sigma\,\kappa\,\mathbf{n}\,\delta(\mathbf{x}-\mathbf{X}),
  \label{eq:tensao}
\end{equation}
espalhada pelo mesmo n\'ucleo. O que \'e novo \'e que $\kappa$ (curvatura) e
$\mathbf{n}$ (normal) v\^em da \emph{geometria} da frente, e portanto da
conectividade dos marcadores --- que o caso r\'igido n\~ao precisava ter.

\paragraph{3. As propriedades saltam.} $\rho$ e $\mu$ mudam atrav\'es da
interface. \textbf{Esta terceira mudan\c{c}a est\'a deliberadamente adiada:} todos
os resultados deste relat\'orio usam propriedades iguais nas duas fases
($\rho=\mu=1$), de modo a isolar a tens\~ao superficial de qualquer efeito de
salto de propriedade. \'E tamb\'em o que torna a compara\c{c}\~ao com o VOF
leg\'itima, como se ver\'a: o caso VOF dispon\'ivel tem exatamente as mesmas
propriedades iguais.

\subsection{Representa\c{c}\~ao da frente e curvatura}

A frente \'e uma polilinha fechada de marcadores \textbf{ordenados}, com a
conectividade impl\'icita na ordem do vetor: o marcador $i$ liga em $i+1$, e o
\'ultimo no primeiro. Isso a distingue do corpo r\'igido, onde os marcadores s\~ao
distribu\'idos por posse euleriana e a ordem da curva \'e destru\'ida --- correto
l\'a, invi\'avel aqui, porque a curvatura precisa de vizinhos.

A curvatura sai do c\'irculo osculador por tr\^es marcadores consecutivos. Para
$\mathbf{p}_{i-1},\mathbf{p}_i,\mathbf{p}_{i+1}$, o circuncentro $\mathbf{c}$
define o vetor curvatura
\begin{equation}
  \kappa\mathbf{n} = \frac{\mathbf{c}-\mathbf{p}_i}{\lVert\mathbf{c}-\mathbf{p}_i\rVert^{2}},
  \label{eq:curvatura}
\end{equation}
de m\'odulo $1/R$ e apontando para o lado c\^oncavo. A forma
\eqref{eq:curvatura} independe da orienta\c{c}\~ao (hor\'aria ou anti-hor\'aria)
da curva, o que elimina uma fonte cl\'assica de erro de sinal.

\paragraph{Uma consequ\^encia a registrar.} Tr\^es pontos colineares d\~ao
$\kappa=0$ --- raio osculador infinito. N\~ao \'e defeito, \'e o m\'etodo; mas
implica que um marcador rec\'em-inserido pela cirurgia, que nasce sobre a corda,
tem curvatura nula at\'e ser advectado para fora da reta.

\subsection{Cirurgia da frente}

Quando a interface estica, os marcadores se afastam e a curva perde
resolu\c{c}\~ao; quando comprime, amontoam-se. A cirurgia reamostra: insere o
ponto m\'edio onde o segmento passa de $2\,\Delta s$, remove marcador onde o
segmento cai abaixo de $\Delta s/2$, com $\Delta s \approx h$.

A remo\c{c}\~ao exige uma guarda que a inser\c{c}\~ao n\~ao exige, e o motivo \'e
geom\'etrico: o ponto m\'edio cai \emph{sobre a corda}, logo a inser\c{c}\~ao \'e
neutra em \'area; j\'a remover o marcador $\mathbf{p}_i$ entre $\mathbf{p}_{i-1}$
e $\mathbf{p}_{i+1}$ descarta o tri\^angulo que eles formam, perdendo \'area
sistematicamente. A remo\c{c}\~ao s\'o acontece, portanto, onde esse tri\^angulo
\'e min\'usculo --- isto \'e, onde a curva \'e quase reta. A
Se\c{c}\~ao~\ref{sec:b1} mostra o que custou descobrir isso por medida.

\paragraph{Mudan\c{c}a de topologia} (gotas coalescendo ou se partindo) \'e o
calcanhar do front-tracking: ao contr\'ario do VOF, n\~ao \'e autom\'atica. Fica
fora do escopo; as fases iniciais assumem topologia fixa.

\section{Implementa\c{c}\~ao}

\subsection{Arquitetura: n\'ucleo puro, adaptador na borda}

A implementa\c{c}\~ao se divide em tr\^es pe\c{c}as, e a divis\~ao foi escolhida
para minimizar risco ao que j\'a funcionava:

\begin{center}
\begin{adjustbox}{max width=\linewidth}
\begin{tabular}{llp{6.6cm}}
\toprule
Pe\c{c}a & Onde & Papel \\
\midrule
n\'ucleo & \texttt{higflow/\allowbreak{}src/\allowbreak{}hig-\allowbreak{}flow-\allowbreak{}front-\allowbreak{}tracking.\allowbreak{}\{c,h\}} &
  frente ordenada, advec\c{c}\~ao, cirurgia, curvatura, \'area. \textbf{N\~ao
  conhece o solver}: compila com \texttt{gcc} simples e \'e test\'avel sozinho \\
adaptador & \texttt{higflow/\allowbreak{}examples-\allowbreak{}common/\allowbreak{}front-\allowbreak{}tracking.\allowbreak{}c} &
  a ponte com o solver: espalha $\sigma\kappa\mathbf{n}$ na malha escalonada e
  interpola $\mathbf{u}$ nos marcadores. Compilado \textbf{por exemplo}, fora da
  biblioteca \\
exemplo & \texttt{higflow/\allowbreak{}example2d\_\allowbreak{}Front\allowbreak{}Tracking} &
  a gota, os par\^ametros e as sondas \\
\bottomrule
\end{tabular}
\end{adjustbox}
\end{center}

Manter o adaptador fora da biblioteca n\~ao \'e detalhe de organiza\c{c}\~ao: a
biblioteca \'e compilada tamb\'em em 3D e \'e ligada por todos os exemplos, de
modo que um erro ali se propaga para casos que nada t\^em com front-tracking. Com
o adaptador por exemplo, quem n\~ao instala uma frente n\~ao referencia
s\'imbolo nenhum, e as su\'ites permanecem intocadas.

\subsection{O reuso que evitou duplicar a parte dif\'icil}

Espalhar e interpolar sobre malha \emph{escalonada} exige achar, para cada
marcador, as facetas no suporte do n\'ucleo --- e essa busca \'e delicada: precisa
ancorar na c\'elula (n\~ao numa faceta, que pode n\~ao existir naquela
dire\c{c}\~ao), cutucar quando o ponto cai exatamente sobre a fronteira de
c\'elula, tratar a conven\c{c}\~ao de identificador negativo em faceta
compartilhada, e detectar suporte atravessando n\'ivel de refinamento. Essa
l\'ogica j\'a existia, testada, dentro do m\'odulo de fronteira imersa r\'igida.

Em vez de reescrev\^e-la, ela foi \textbf{exposta} (\texttt{fi\_suporte\_facetas}),
num invólucro que n\~ao altera o comportamento do corpo r\'igido. O adaptador de
front-tracking a consome. \'E a diferen\c{c}a entre reusar c\'odigo verificado e
criar uma segunda c\'opia para divergir com o tempo.

\subsection{Acoplamento: onde a for\c{c}a entra no passo}

O gancho do front-tracking roda \textbf{antes do preditor}, no mesmo ponto em que
a fronteira imersa r\'igida injeta sua for\c{c}a, e depois que o termo fonte de
faceta foi escrito. Nesse instante o campo \texttt{dpu} guarda $\mathbf{u}^{n}$ ---
a velocidade final do passo anterior, \emph{j\'a projetada} e portanto
discretamente livre de diverg\^encia. A ordem dentro do gancho \'e:

\begin{enumerate}
  \item move a frente com $\mathbf{u}^{n}$ (advec\c{c}\~ao);
  \item cirurgia;
  \item espalha $\sigma\kappa\mathbf{n}$ nas posi\c{c}\~oes novas, somando no campo
    de fonte que a equa\c{c}\~ao de momento l\^e.
\end{enumerate}

Como n\~ao existe $\mathbf{u}$ em $t+\Delta t/2$, o campo fica congelado em
$\mathbf{u}^{n}$ dentro do passo: o RK2 de ponto m\'edio usado na advec\c{c}\~ao
vira avalia\c{c}\~ao de ponto m\'edio no \emph{espa\c{c}o}, e o esquema \'e de
primeira ordem no tempo --- coerente com o restante do acoplamento expl\'icito.

\paragraph{Limita\c{c}\~ao atual: serial.} O espalhamento n\~ao faz a
redu\c{c}\~ao das facetas de franja para os donos, que em $n_p>1$ seria
necess\'aria (o m\'odulo r\'igido a faz com \texttt{PetscSFReduce}). Todos os
resultados deste relat\'orio s\~ao $n_p=1$, onde n\~ao h\'a franja.

\section{Verifica\c{c}\~ao}
\label{sec:verificacao}

Cada pe\c{c}a foi verificada contra um or\'aculo \emph{exato} antes de a
pr\'oxima ser constru\'ida. Os arreios est\~ao em
\texttt{higflow/tests/front-tracking} e devolvem c\'odigo de sa\'ida: s\~ao
port\~oes, n\~ao impress\~oes.

\subsection{B1: advec\c{c}\~ao cinem\'atica reversivel}
\label{sec:b1}

O teste de Rider--Kothe~\cite{rider1998}: um c\'irculo advectado no campo de
v\'ortice \'unico
\begin{align}
  u &= \phantom{-}\sin^{2}(\pi x)\sin(2\pi y)\cos(\pi t/T),\\
  v &= -\sin(2\pi x)\sin^{2}(\pi y)\cos(\pi t/T),
\end{align}
que \'e livre de diverg\^encia e se inverte em $t=T/2$, de modo que a frente deve
voltar \`a forma inicial em $t=T$. Dois or\'aculos: a \emph{conserva\c{c}\~ao} de
\'area ao longo de toda a corrida, e a \emph{revers\~ao} em $t=T$.

\pgfplotstabletypeset[
  columns/nmarc/.style={column name={$N$}, int detect},
  columns/dt/.style={column name={$\Delta t$}, fixed, precision=3},
  columns/conssem/.style={column name={cons.\ sem cir.}, sci, sci zerofill, precision=2},
  columns/revsem/.style={column name={rev.\ sem cir.}, sci, sci zerofill, precision=2},
  columns/conscom/.style={column name={cons.\ com cir.}, sci, sci zerofill, precision=2},
  columns/revcom/.style={column name={rev.\ com cir.}, sci, sci zerofill, precision=2},
  columns/nfinal/.style={column name={$N$ final}, int detect},
  every head row/.style={before row=\toprule, after row=\midrule},
  every last row/.style={after row=\bottomrule},
]{\dat{b1-rider-kothe.dat}}

A advec\c{c}\~ao \'e de segunda ordem: sem cirurgia, a revers\~ao cai
\num{3.2e-4} $\to$ \num{4.0e-5} $\to$ \num{5.0e-6}. Com cirurgia, a
conserva\c{c}\~ao de \'area fica \emph{uma ordem de grandeza melhor} sob
esticamento, porque manter $\Delta s \approx h$ preserva a resolu\c{c}\~ao do
filamento; em troca, a revers\~ao perfeita se perde, o que \'e o compromisso
correto.


\begin{figure}[htbp]
\centering
\begin{tikzpicture}
\begin{axis}[
  width=11cm, height=8cm, axis equal image,
  xlabel={$x$}, ylabel={$y$}, grid=major,
  xmin=0.1, xmax=0.9, ymin=0.3, ymax=1.0,
  legend pos=south west, legend cell align=left, legend style={font=\small},
]
\addplot[thick]              table[x=x,y=y] {\dat{formas/b1_0.dat}};
\addlegendentry{$t=0$ (c\'irculo)}
\addplot[thin, dashed]       table[x=x,y=y] {\dat{formas/b1_2.dat}};
\addlegendentry{$t=T/2$ (filamento)}
\addplot[thick, densely dotted] table[x=x,y=y] {\dat{formas/b1_4.dat}};
\addlegendentry{$t=T$ (retorno)}
\end{axis}
\end{tikzpicture}
\caption{Teste revers\'ivel de Rider--Kothe. A frente estica em filamento at\'e
$t=T/2$, quando o campo se inverte, e retorna \`a forma inicial em $t=T$. A
cirurgia acompanha o esticamento inserindo marcadores (256 $\to$ 680) e os
remove no retorno (654). A \'area fechada varia \num{6e-5} ao longo de todo o
percurso.}
\label{fig:b1}
\end{figure}

\paragraph{O que custou uma itera\c{c}\~ao.} A primeira vers\~ao da cirurgia
removia marcador de segmento curto \emph{sem guarda de curvatura}, e cortava
canto: o erro de \'area passou a ser dominado pela pr\'opria cirurgia, com a
assinatura inconfund\'ivel de o erro m\'aximo coincidir com o erro final. A
guarda descrita na Se\c{c}\~ao~2.3 corrigiu, e o teste mostrou a melhora.

\subsection{Curvatura}

Contra dois or\'aculos anal\'iticos: c\'irculo ($\kappa=1/R$ em todo marcador) e
elipse ($\kappa(t) = ab/(a^{2}\sin^{2}t + b^{2}\cos^{2}t)^{3/2}$).

\pgfplotstabletypeset[
  columns/n/.style={column name={$N$ (c\'irculo)}, int detect},
  columns/circulo/.style={column name={erro rel.\ c\'irculo}, sci, sci zerofill, precision=2},
  columns/nelipse/.style={column name={$N$ (elipse)}, int detect},
  columns/elipse/.style={column name={erro rel.\ elipse}, sci, sci zerofill, precision=2},
  every head row/.style={before row=\toprule, after row=\midrule},
  every last row/.style={after row=\bottomrule},
]{\dat{curvatura.dat}}

No c\'irculo a curvatura \'e \textbf{exata} em precis\~ao de m\'aquina --- tr\^es
pontos sobre um c\'irculo definem esse mesmo c\'irculo. Na elipse converge em
segunda ordem. Esse resultado \'e o que autoriza, mais adiante, usar
$\sigma/R$ como refer\^encia exata para o salto de press\~ao: a curvatura que a
for\c{c}a \eqref{eq:tensao} aplica numa gota circular \'e $1/R$ sem erro de
discretiza\c{c}\~ao.

\subsection{Espalhamento da for\c{c}a}

Dois or\'aculos sobre uma gota em repouso, ambos em precis\~ao de m\'aquina:
\emph{conserva\c{c}\~ao} --- $\sum_{\text{malha}} \mathbf{F}\,h^{2}$ deve igualar
$\sigma\sum_k \kappa\mathbf{n}_k\,\mathrm{d}s_k$, porque o n\'ucleo de Roma soma
1 (parti\c{c}\~ao da unidade) --- e \emph{equil\'ibrio}, $\oint \kappa\mathbf{n}\,
\mathrm{d}s \to 0$ para curva fechada, que \'e a for\c{c}a l\'iquida nula exigida
pela lei de Laplace.

\pgfplotstabletypeset[
  columns/nmarc/.style={column name={$N$}, int detect},
  columns/conservx/.style={column name={conserv.\ $x$}, sci, sci zerofill, precision=2},
  columns/conservy/.style={column name={conserv.\ $y$}, sci, sci zerofill, precision=2},
  columns/equilib/.style={column name={equil\'ibrio $\lvert\mathbf{F}\rvert$}, sci, sci zerofill, precision=2},
  every head row/.style={before row=\toprule, after row=\midrule},
  every last row/.style={after row=\bottomrule},
]{\dat{tensao-conservacao.dat}}

\paragraph{Honestidade sobre o segundo or\'aculo.} O equil\'ibrio dar
m\'aquina-zero vem da \emph{simetria do c\'irculo}; uma curva fechada
assim\'etrica daria um valor finito convergindo. Como o c\'irculo \'e exatamente
o caso de Laplace, m\'aquina-zero \'e a resposta certa aqui --- mas o teste forte
desta subse\c{c}\~ao \'e a conserva\c{c}\~ao, que n\~ao depende de simetria.

\subsection{Advec\c{c}\~ao acoplada}

Com o solver no circuito, a frente passa a se mover com a velocidade
interpolada da malha. A gota est\'atica serve de or\'aculo: ela deve ficar parada,
circular e com \'area preservada, j\'a que o campo projetado \'e livre de
diverg\^encia.

\pgfplotstabletypeset[
  columns/passo/.style={column name={passo}, int detect},
  columns/n/.style={column name={$N$}, int detect},
  columns/area/.style={column name={\'area}, fixed, precision=8},
  columns/driftrel/.style={column name={$\Delta A/A$}, sci, sci zerofill, precision=2},
  columns/deriva/.style={column name={deriva centr.}, sci, sci zerofill, precision=2},
  columns/circ/.style={column name={circularidade}, fixed, precision=6},
  columns/unidade/.style={column name={part.\ unidade}, sci, sci zerofill, precision=2},
  columns/umarc/.style={column name={$\max\lvert u_{\text{marc}}\rvert$}, sci, sci zerofill, precision=2},
  every head row/.style={before row=\toprule, after row=\midrule},
  every last row/.style={after row=\bottomrule},
]{\dat{adveccao-acoplada.dat}}

A \emph{parti\c{c}\~ao da unidade} em \num{6.9e-15} \'e o or\'aculo forte: os
pesos somam 1 em precis\~ao de m\'aquina, o que prova que um campo uniforme seria
interpolado exatamente e pega suporte incompleto --- que daria velocidade pequena
demais, em sil\^encio.

\paragraph{A checagem que evitou um falso verde.} Uma gota que fica parada \'e
\emph{exatamente} o que se veria se a interpola\c{c}\~ao devolvesse zero em
sil\^encio: \'area perfeita, deriva nula, circularidade constante. Da\'i a
\'ultima coluna. Ela vale \num{2.9e-4} --- n\~ao nula --- e vale
\emph{exatamente zero no passo 0}, quando a condi\c{c}\~ao inicial \'e
$\mathbf{u}=0$. O n\'umero acompanha o campo, logo a advec\c{c}\~ao est\'a de
fato operando.

\paragraph{O que estes testes n\~ao exercitam.} A gota est\'atica n\~ao deforma:
$N$ permanece constante e \textbf{a cirurgia nunca dispara} no regime acoplado.
O caveat de $\kappa=0$ em marcador rec\'em-inserido segue, portanto, sem
verifica\c{c}\~ao com o solver no circuito. Ambos entram em cena em B3 e B4.

\section{Compara\c{c}\~ao com o VOF}
\label{sec:comparacao}

\subsection{Por que a compara\c{c}\~ao \'e leg\'itima}

Comparar dois m\'etodos s\'o informa se eles resolverem o mesmo problema. O
HiGFlow j\'a possu\'ia um caso VOF que, examinado, \emph{\'e} o teste de gota
est\'atica de Laplace, sem nenhuma adapta\c{c}\~ao: gota circular de raio
$R=1/6$ centrada em $(0{,}5;\,0{,}5)$, dom\'inio $[0,1]^2$, duas fases
newtonianas, tens\~ao superficial ligada e velocidade inicial nula.

Trocando em mi\'udos, os dois lados coincidem em tudo que importa:

\begin{center}
\begin{tabular}{lll}
\toprule
 & front-tracking & VOF \\
\midrule
dom\'inio            & \multicolumn{2}{c}{$[0,1]^2$, malha $61\times61$, $h=1/61$} \\
gota                 & \multicolumn{2}{c}{c\'irculo $R=1/6$ em $(0{,}5;\,0{,}5)$} \\
c\'elulas por raio   & \multicolumn{2}{c}{$10{,}2$} \\
propriedades         & \multicolumn{2}{c}{$\rho=\mu=1$ nas duas fases} \\
$\sigma$ efetivo     & \multicolumn{2}{c}{$1$} \\
$Re$, $\Delta t$, passos & \multicolumn{2}{c}{$1$; $10^{-3}$; $101$} \\
salto esperado       & \multicolumn{2}{c}{$\Delta p = \sigma/R = 6$} \\
\midrule
interface            & polilinha de 64 marcadores & fra\c{c}\~ao volum\'etrica \\
tens\~ao superficial & $\sigma\kappa\mathbf{n}$ espalhada (Roma) & CSF~\cite{brackbill1992} \\
\bottomrule
\end{tabular}
\end{center}

O $\sigma$ efetivo merece uma palavra, porque n\~ao \'e \'obvio que coincida. No
caminho multif\'asico, a for\c{c}a interfacial entra na equa\c{c}\~ao de momento
como $\mathbf{IF}/(We\,\rho)$ com $We = Ca\cdot Re$. Com $Ca=Re=1$ e $\rho=1$,
resulta $We=1$ e portanto $\sigma_{\text{ef}}=1$ --- exatamente o $\sigma$ que o
front-tracking usa explicitamente. A coincid\^encia \'e feliz e foi verificada no
c\'odigo, n\~ao suposta.

Al\'em disso: mesmo solver, mesmo passo no tempo, mesmas sondas de press\~ao nos
mesmos pontos, e as correntes par\'asitas lidas do \emph{mesmo} bloco de
impress\~ao do solver nos dois casos.

\paragraph{Um confundidor que precisou ser removido.} O caso VOF, como
distribu\'ido, imp\~oe na parede superior uma velocidade que \emph{cresce no
tempo}, $8\,(1+\tanh(8t-4))\,x^{2}(1-x)^{2}$. Com ela ligada, a velocidade medida
n\~ao \'e corrente par\'asita: \'e par\'asita mais escoamento for\c{c}ado, e
compar\'a-la com a gota est\'atica do front-tracking mediria coisas diferentes.
O for\c{c}amento foi zerado para a compara\c{c}\~ao, atr\'as de um seletor que
preserva o comportamento padr\~ao do exemplo.

\subsection{Resultados}

\pgfplotstabletypeset[
  string type,
  columns/metrica/.style={column name={m\'etrica}},
  columns/ft/.style={column name={FT espalhado}},
  columns/ftbal/.style={column name={FT balanceado}},
  columns/vof/.style={column name={VOF}},
  columns/melhor/.style={column name={melhor}},
  every head row/.style={before row=\toprule, after row=\midrule},
  every last row/.style={after row=\bottomrule},
]{\dat{comparacao-laplace.dat}}

Em valores brutos:

\pgfplotstabletypeset[
  string type,
  columns/grandeza/.style={column name={grandeza}},
  columns/ft/.style={column name={FT espalhado}},
  columns/ftbal/.style={column name={FT balanceado}},
  columns/vof/.style={column name={VOF}},
  columns/exato/.style={column name={exato}},
  every head row/.style={before row=\toprule, after row=\midrule},
  every last row/.style={after row=\bottomrule},
]{\dat{comparacao-bruto.dat}}

\subsection{Leitura dos tr\^es resultados}

\paragraph{Salto de press\~ao: o front-tracking ganha, seis vezes.}
$\Delta p = 5{,}9954$ contra $6{,}0284$, para o valor exato $6$. \'E o resultado
esperado, e a raz\~ao \'e direta: a curvatura do front-tracking sai da geometria
expl\'icita da frente e \'e exata para um c\'irculo (Se\c{c}\~ao~\ref{sec:verificacao}),
enquanto a do VOF vem de derivadas de um campo de fra\c{c}\~ao suavizado --- o
ponto historicamente fr\'agil do m\'etodo.

\paragraph{Conserva\c{c}\~ao de massa: o VOF ganha, por nove ordens de grandeza.}
A deriva no tempo do VOF \'e \num{2.9e-15}, isto \'e, zero de m\'aquina: ele
advecta um escalar conservado, e conservar \'e propriedade de constru\c{c}\~ao.
O front-tracking deriva \num{1.3e-5} em 100 passos. Tamb\'em esperado, e \'e o
outro lado da mesma moeda.

Aqui vale registrar um erro de medi\c{c}\~ao cometido e corrigido no caminho:
a primeira compara\c{c}\~ao punha os dois lados em \num{1.3e-5} e parecia empate.
Era artefato de medir coisas diferentes --- para o VOF, desvio contra a \'area
anal\'itica; para o front-tracking, deriva contra a \'area inicial. Separadas as
duas quantidades, a \num{1.3e-5} do VOF revelou-se erro de \emph{inicializa\c{c}\~ao}
(a discretiza\c{c}\~ao do c\'irculo em $t=0$) e a deriva verdadeira, zero de
m\'aquina.

\paragraph{Correntes par\'asitas: na forma espalhada o VOF ganha, dez vezes ---
e isso contraria a expectativa.} \'E o resultado que motivou a
Se\c{c}\~ao~\ref{sec:balanco}, onde a causa \'e identificada e removida: a
coluna ``FT balanceado'' das tabelas acima j\'a traz o desfecho, e vem de l\'a.
Na forma espalhada, as par\'asitas \emph{estacionam} num patamar; na balanceada,
\emph{decaem} at\'e o piso de impress\~ao do solver --- a diferen\'ca entre um
for\c{c}amento esp\'urio persistente e a aus\^encia dele.

\begin{center}
\begin{tikzpicture}
\begin{semilogyaxis}[
  width=12cm, height=6.4cm,
  xlabel={passo}, ylabel={$\max\lvert u\rvert$},
  legend pos=north east, grid=major,
  legend cell align=left,
]
\addplot table[x=passo, y=ft]    {\dat{parasitas-serie.dat}};
\addlegendentry{FT espalhado}
\addplot table[x=passo, y=vof]   {\dat{parasitas-serie.dat}};
\addlegendentry{VOF}
\addplot table[x=passo, y=ftbal] {\dat{parasitas-serie.dat}};
\addlegendentry{FT balanceado}
\end{semilogyaxis}
\end{tikzpicture}
\end{center}

\subsection{Erro de forma na inicializa\c{c}\~ao}

A linha da inicializa\c{c}\~ao merece ressalva. O d\'eficit de \num{1.6e-3} do
front-tracking \'e o d\'eficit de \'area de um pol\'igono de 64 lados inscrito no
c\'irculo --- artefato de \emph{como se inicializa}, n\~ao propriedade do
m\'etodo. Ele melhora com mais marcadores, e desapareceria ao inscrever o
pol\'igono com a \'area correta. O VOF, que integra a fra\c{c}\~ao com
subdivis\~ao $16\times16$ por c\'elula, parte muito mais pr\'oximo. A compara\c{c}\~ao
nesta linha \'e entre escolhas de implementa\c{c}\~ao, e n\~ao entre m\'etodos.

\section{A hip\'otese do balan\c{c}o, testada}
\label{sec:balanco}

\subsection{A hip\'otese}

A compara\c{c}\~ao deixou um resultado desconfort\'avel: o front-tracking obt\'em
a curvatura \emph{melhor} --- exata no c\'irculo --- e ainda assim produz
correntes par\'asitas dez vezes maiores. Os dois fatos s\'o convivem se o
problema n\~ao estiver na curvatura.

A hip\'otese \'e que ele est\'a no \textbf{balan\c{c}o discreto}. Para a gota em
equil\'ibrio o cont\'inuo exige
\begin{equation}
  \nabla p = \sigma\kappa\mathbf{n}\,\delta_S,
  \label{eq:equil}
\end{equation}
e a velocidade par\'asita mede exatamente o quanto os \emph{operadores
discretos} dos dois lados falham em satisfazer \eqref{eq:equil}. N\~ao basta que
cada lado esteja individualmente correto: \'e preciso que a for\c{c}a discreta
\textbf{perten\c{c}a \`a imagem do gradiente discreto} que a proje\c{c}\~ao
inverte. \'E a propriedade de \emph{balanced-force} de Francois et
al.~\cite{francois2006}.

No VOF com CSF a for\c{c}a \'e $\sigma\kappa\nabla f$; numa gota circular $\kappa$
\'e constante, de sorte que $\sigma\kappa\nabla f = \nabla(\sigma\kappa f)$ ---
a for\c{c}a \emph{\'e} um gradiente discreto, e pode ser cancelada. Na forma
espalhada do front-tracking, a for\c{c}a sai do n\'ucleo de Roma sobre facetas
escalonadas, e esse operador n\~ao \'e o gradiente discreto de escalar nenhum que
a proje\c{c}\~ao produza.

\subsection{A reformula\c{c}\~ao}

Para testar, a for\c{c}a foi remontada como
\begin{equation}
  \mathbf{F} = \sigma\,\kappa\,\nabla H,
  \label{eq:balanceada}
\end{equation}
com duas exig\^encias:

\begin{enumerate}
  \item $H$ \'e a \textbf{fra\c{c}\~ao de \'area da gota por c\'elula}, calculada
    da \emph{geometria} da frente por recorte de pol\'igono
    (Sutherland--Hodgman) seguido da f\'ormula do la\c{c}o. \'E o an\'alogo exato
    do campo de cor do VOF, s\'o que derivado da interface em vez de advectado.
  \item $\nabla$ \'e o \textbf{mesmo} operador discreto que
    \texttt{higflow\_\allowbreak{}final\_\allowbreak{}velocity} usa para corrigir a velocidade pela
    press\~ao --- a mesma amostragem das duas c\'elulas vizinhas e a mesma
    diferen\c{c}a dividida. Esta \'e a exig\^encia que faz o m\'etodo
    ``balanceado''; qualquer outro gradiente, ainda que consistente, n\~ao
    cancelaria.
\end{enumerate}

Com $\kappa$ constante --- o caso do c\'irculo --- \eqref{eq:balanceada} \'e
exatamente $\nabla(\sigma\kappa H)$, e a press\~ao pode cancel\'a-la termo a
termo.

\paragraph{A indicadora foi verificada antes de ser usada.} Somar
$H$ vezes a \'area da c\'elula sobre uma grade que cobre a frente tem de devolver
a \'area do pol\'igono, exatamente, porque as c\'elulas particionam o plano.
Medido sobre quatro frentes e quatro grades: pior erro \num{3.7e-15}. C\'elula
inteiramente dentro devolve a \'area da c\'elula (\num{4.0e-14}); inteiramente
fora devolve zero exato.

\subsection{Resultado}

Mesmo caso, mesma malha, mesma gota, mesmo tudo --- trocando apenas a
formula\c{c}\~ao da for\c{c}a:

\begin{center}
\begin{adjustbox}{max width=\linewidth}
\begin{tabular}{lrrr}
\toprule
 & FT espalhado & FT balanceado & fator \\
\midrule
erro relativo de $\Delta p$        & \num{7.67e-04} & \num{1.7e-07}  & $4500\times$ \\
correntes par\'asitas (final)      & \num{5.06e-04} & $<\num{5e-11}$ & $>10^{4}\times$ \\
deriva de \'area                   & \num{1.29e-05} & \num{1.9e-15}  & $7\times10^{9}$ \\
deriva do centroide                & \num{3.4e-07}  & \num{1.9e-14}  & $2\times10^{7}$ \\
\bottomrule
\end{tabular}
\end{adjustbox}
\end{center}

\textbf{A hip\'otese se confirma, e nos dois sentidos que ela previa:} as
correntes par\'asitas caem mais de quatro ordens de grandeza \emph{e} o salto de
press\~ao \emph{melhora} --- n\~ao piora. Fosse a curvatura o problema, n\~ao
haveria como as duas coisas andarem juntas.

\paragraph{A assinatura temporal \'e ainda mais clara que os n\'umeros finais.}
Na forma espalhada as par\'asitas sobem e \emph{estacionam} num patamar de
\num{5e-4}: h\'a um for\c{c}amento esp\'urio persistente, que a viscosidade
equilibra mas n\~ao elimina. Na balanceada elas decaem monotonicamente ---
\num{2.4e-7}, \num{1.6e-9}, \num{2e-10} --- at\'e sumirem sob a precis\~ao de
impress\~ao do solver. N\~ao h\'a for\c{c}amento esp\'urio a equilibrar; o que se
v\^e \'e um transiente inicial sendo amortecido.

\paragraph{E a for\c{c}a n\~ao sumiu.} Velocidade nula \'e tamb\'em o que se veria
se a for\c{c}a simplesmente n\~ao estivesse sendo aplicada --- o mesmo modo de
falso verde que j\'a apareceu na verifica\c{c}\~ao da advec\c{c}\~ao. O que
descarta essa leitura \'e o pr\'oprio salto de press\~ao: $\Delta p = 6{,}000001$
contra o exato $6$. Sem for\c{c}a n\~ao haveria salto nenhum. As duas medidas
juntas --- press\~ao certa \emph{e} velocidade nula --- s\~ao a assinatura do
equil\'ibrio, e nenhuma delas sozinha bastaria.

\subsection{Consequ\^encia para a compara\c{c}\~ao}

Com a for\c{c}a balanceada, o front-tracking passa a superar o VOF nas tr\^es
m\'etricas medidas, e n\~ao mais em uma s\'o:

\begin{center}
\begin{adjustbox}{max width=\linewidth}
\begin{tabular}{lrrl}
\toprule
 & FT balanceado & VOF & vantagem \\
\midrule
erro relativo de $\Delta p$   & \num{1.7e-07}  & \num{4.73e-03} & $2{,}8\times10^{4}$ \\
correntes par\'asitas         & $<\num{5e-11}$ & \num{5.23e-05} & $>10^{6}$ \\
deriva de massa               & \num{1.9e-15}  & \num{2.9e-15}  & equivalentes \\
\bottomrule
\end{tabular}
\end{adjustbox}
\end{center}

A conserva\c{c}\~ao de massa deixou de ser desvantagem do front-tracking, mas por
uma raz\~ao que convém explicitar para n\~ao superestimar o resultado: a \'area
deixa de derivar porque o escoamento \emph{\'e} est\'atico, e marcadores parados
n\~ao perdem \'area. Num escoamento de verdade a deriva volta, e a
conserva\c{c}\~ao por constru\c{c}\~ao do VOF --- em malha fixa, ver a
Se\c{c}\~ao~\ref{sec:massa-amr} --- continua sendo uma vantagem
estrutural que este teste n\~ao mede.

\paragraph{O que o teste custou, em c\'odigo.} Uma fun\c{c}\~ao de geometria no
n\'ucleo (\texttt{ft\_area\_na\_caixa}, com seu arreio) e uma rotina de montagem
no adaptador. A formula\c{c}\~ao antiga continua sendo a padr\~ao; a balanceada
entra por seletor, de modo que os dois caminhos coexistem e podem ser comparados
a qualquer momento --- que foi, afinal, como esta se\c{c}\~ao p\^ode ser escrita.

\section{Gota el\'iptica: curvatura vari\'avel}
\label{sec:elipse}

\subsection{O teste que faltava}

A Se\c{c}\~ao~\ref{sec:balanco} fechou com uma ressalva expl\'icita: o balan\c{c}o
\'e \emph{exato} apenas quando $\kappa$ \'e constante, pois a\'i
$\sigma\kappa\nabla H = \nabla(\sigma\kappa H)$. Com curvatura vari\'avel sobra
um res\'iduo, e a gota circular n\~ao o mede. Esta se\c{c}\~ao mede.

O caso: elipse de semieixos $a=0{,}2357$ e $b=0{,}1179$ --- raz\~ao $2{:}1$ e
\'area \emph{id\^entica} \`a do c\'irculo $R=1/6$ das se\c{c}\~oes anteriores.
A gota n\~ao est\'a em equil\'ibrio; a tens\~ao superficial a relaxa para o
c\'irculo de mesma \'area, de modo que o raio de equil\'ibrio \'e
$R_{\text{eq}}=\sqrt{ab}=1/6$ e o salto final esperado volta a ser
$\Delta p = 6$. Mesma malha, mesmas propriedades, 1200 passos ($t=1{,}2$, cerca
de quatro tempos de relaxa\c{c}\~ao).

\'E tamb\'em o primeiro caso em que a frente \textbf{deforma}, e portanto o
primeiro em que a \textbf{cirurgia dispara} com o solver no circuito: o n\'umero
de marcadores cai de 108 para 102--104 conforme a elipse se contrai no eixo
maior.

\subsection{O defeito que o c\'irculo escondia}

Na primeira corrida, a formula\c{c}\~ao balanceada \textbf{falhou}. A
circularidade sobe normalmente at\'e o passo~900, chega a $0{,}9887$ --- e ent\~ao
\emph{reverte}, caindo a $0{,}9617$, com a velocidade crescendo. O salto de
press\~ao final sai $10{,}99$ contra os $6$ esperados: \SI{83}{\percent} de erro.

\begin{center}
\begin{tikzpicture}
\begin{axis}[
  width=12cm, height=6.4cm,
  xlabel={passo}, ylabel={circularidade $4\pi A/P^{2}$},
  legend pos=south east, grid=major, legend cell align=left,
  ymin=0.83, ymax=1.01,
]
\addplot table[x=passo, y=espalhado]  {\dat{elipse-circularidade.dat}};
\addlegendentry{espalhado}
\addplot table[x=passo, y=balsuave]   {\dat{elipse-circularidade.dat}};
\addlegendentry{balanceado, $\kappa$ suavizado}
\addplot table[x=passo, y=balvizinho] {\dat{elipse-circularidade.dat}};
\addlegendentry{balanceado, $\kappa$ do vizinho}
\end{axis}
\end{tikzpicture}
\end{center}

O instante da falha \'e a pista: ela aparece justamente quando a gota fica quase
circular, isto \'e, quando a for\c{c}a f\'isica que a movia j\'a decaiu. Enquanto
havia um sinal grande, o artefato ficava submerso; ao ficar sozinho, dominou.

A causa estava numa simplifica\c{c}\~ao que a pr\'opria
Se\c{c}\~ao~\ref{sec:discussao} havia registrado como limita\c{c}\~ao: a
curvatura na faceta era a do \textbf{marcador mais pr\'oximo}. No c\'irculo isso
\'e exato, porque $\kappa$ \'e constante. Na elipse \'e um campo
\emph{descont\'inuo}: salta quando a atribui\c{c}\~ao troca de marcador. As
descontinuidades produzem saltos de for\c{c}a que movem a frente, o que muda a
atribui\c{c}\~ao, o que muda a for\c{c}a --- realimenta\c{c}\~ao.

\subsection{Conserto dirigido, e a confirma\c{c}\~ao do diagn\'ostico}

Trocou-se apenas isso: a curvatura na faceta passou a ser a \textbf{m\'edia
ponderada} das curvaturas dos marcadores vizinhos, com o n\'ucleo de Roma de
largura $h$. Nada mais mudou.

Com $\kappa$ suavizado a instabilidade \textbf{desaparece}, e a trajet\'oria
passa a coincidir com a da formula\c{c}\~ao espalhada at\'e a quarta casa:
$0{,}922213$ contra $0{,}922220$ no passo~200; $0{,}998477$ contra $0{,}998491$
no passo~1150. Duas formula\c{c}\~oes de for\c{c}a diferentes produzindo a mesma
din\^amica \'e evid\^encia forte de que ambas est\~ao certas.

\paragraph{Sem regress\~ao no c\'irculo.} A curvatura \'e constante ali, e a
m\'edia ponderada de um valor constante devolve esse valor exatamente: o
resultado da Se\c{c}\~ao~\ref{sec:balanco} se mant\'em --- $\Delta p = 6{,}000001$
e $\max\lvert u\rvert < \num{5e-11}$.

\subsection{Resultados, e o que eles corrigem}

\pgfplotstabletypeset[
  string type,
  columns/grandeza/.style={column name={grandeza}},
  columns/espalhado/.style={column name={espalhado}},
  columns/balvizinho/.style={column name={bal., $\kappa$ vizinho}},
  columns/balsuave/.style={column name={bal., $\kappa$ suave}},
  columns/vof/.style={column name={VOF}},
  every head row/.style={before row=\toprule, after row=\midrule},
  every last row/.style={after row=\bottomrule},
]{\dat{elipse-resultados.dat}}

Tr\^es leituras, e duas delas corrigem impress\~oes deixadas pelas se\c{c}\~oes
anteriores.

\paragraph{1. Na elipse, o balan\c{c}o n\~ao traz vantagem.} Espalhado
\SI{1.27}{\percent}, balanceado \SI{1.52}{\percent}, VOF \SI{1.19}{\percent} ---
os tr\^es dentro de \SI{0.3}{\percent} um do outro. O ganho espetacular da
Se\c{c}\~ao~\ref{sec:balanco} era \emph{espec\'ifico do caso est\'atico com
$\kappa$ constante}, onde o balan\c{c}o exato \'e alcan\c{c}\'avel. Sob
movimento real e curvatura vari\'avel, as duas formula\c{c}\~oes s\~ao
equivalentes dentro da medida.

\paragraph{2. A vantagem estrutural do VOF em massa aparece agora.} A
Se\c{c}\~ao~\ref{sec:balanco} advertiu que a conserva\c{c}\~ao perfeita do
front-tracking balanceado vinha de o escoamento \emph{ser} est\'atico, e que ``em
escoamento de verdade a deriva volta''. Voltou: \num{1.1e-4} e \num{2.3e-4}
contra \num{2.5e-13} do VOF. Este \'e o primeiro teste com movimento real, e ele
confirma a advert\^encia.

\paragraph{3. Os tr\^es m\'etodos concordam na f\'isica.} Os saltos finais ficam
em $5{,}909$, $5{,}924$ e $5{,}929$ para um valor exato de $6$. O d\'eficit comum
de \SI{1.2}{\percent} n\~ao \'e de m\'etodo: \'e de relaxa\c{c}\~ao incompleta ---
em $t=1{,}2$ a circularidade ainda \'e $0{,}9985$, e resta velocidade da ordem de
\num{9e-3}, que contribui press\~ao din\^amica. Que tr\^es discretiza\c{c}\~oes
independentes errem \emph{junto}, e pelo mesmo tanto, \'e o que se espera de um
erro f\'isico compartilhado, e n\~ao de defeito em qualquer uma delas.

\subsection{A li\c{c}\~ao de m\'etodo}

O resultado da gota circular era verdadeiro e continua v\'alido; o que faltava
era saber \emph{sobre o que} ele falava. Um caso em que uma quantidade \'e
constante n\~ao distingue uma implementa\c{c}\~ao que trata bem essa quantidade
de outra que a trata de qualquer jeito --- e aqui a diferen\c{c}a era entre
funcionar e ir \`a instabilidade.

A limita\c{c}\~ao estava registrada por escrito antes de ser testada, o que
tornou o diagn\'ostico r\'apido: bastou olhar para o \'unico ingrediente que
mudava de comportamento entre c\'irculo e elipse. Registrar uma limita\c{c}\~ao
conhecida no momento em que ela \'e aceita paga-se quando ela cobra.

\section{Discuss\~ao}
\label{sec:discussao}

\subsection{O que os dois m\'etodos oferecem, pelo que foi medido}

\begin{center}
\begin{tabular}{p{4.6cm}p{4.6cm}p{4.6cm}}
\toprule
 & front-tracking (balanceado) & VOF \\
\midrule
curvatura / salto de press\~ao & exata no c\'irculo; erro \num{1.7e-7}
  & derivada de campo suavizado; erro \num{4.7e-3} \\
correntes par\'asitas & $<\num{5e-11}$, decaindo & \num{5.2e-5}, decaindo \\
massa & conservada enquanto o escoamento \'e est\'atico; deriva com o movimento
  & conservada por constru\c{c}\~ao, sempre \\
topologia (coalesc\^encia, quebra) & manual, fora do escopo & autom\'atica \\
custo de estrutura & marcadores ordenados, cirurgia, indicadora por c\'elula
  & campo euleriano a mais \\
\bottomrule
\end{tabular}
\end{center}

A leitura n\~ao \'e ``um vence''. Nas m\'etricas deste teste o front-tracking
balanceado sai \`a frente, mas o teste \'e uma gota est\'atica: n\~ao exercita
topologia, n\~ao exercita deforma\c{c}\~ao grande, e n\~ao exercita a
conserva\c{c}\~ao de massa sob movimento --- que \'e justamente onde o VOF tem
vantagem \emph{estrutural}, e n\~ao circunstancial. Escoamento em que a massa da
fase precisa ser conservada ao longo de muitos tempos, ou em que gotas coalescem
e se partem, continua favorecendo o VOF. Com os dois no mesmo sistema, a escolha
passa a ser informada por medida em vez de por prefer\^encia.

\subsection{A compara\c{c}\~ao como verifica\c{c}\~ao m\'utua}

Vale sublinhar um uso da compara\c{c}\~ao que n\~ao \'e o de escolher m\'etodo.
Como a dire\c{c}\~ao de v\'arios resultados \'e conhecida \emph{a priori} pela
literatura, eles funcionam como or\'aculo um para o outro: se o VOF tivesse
sa\'ido melhor em curvatura, ou se o front-tracking espalhado tivesse conservado
massa exatamente, o resultado apontaria defeito de implementa\c{c}\~ao, n\~ao
descoberta sobre os m\'etodos.

E foi exatamente assim que o defeito apareceu. O \'unico resultado que
\emph{contrariou} a expectativa --- correntes par\'asitas maiores no
front-tracking --- n\~ao era propriedade do m\'etodo: era uma deficiência da
implementa\c{c}\~ao, localizada pela discrep\^ancia e removida na
Se\c{c}\~ao~\ref{sec:balanco}. Sem o VOF ao lado, com o mesmo problema e as
mesmas sondas, n\~ao haveria com o que estranhar \num{5e-4}: pareceria apenas
``pequeno''.

\subsection{Limita\c{c}\~oes}

Para que o alcance dos n\'umeros n\~ao seja superestimado:

\begin{itemize}
  \item \textbf{Uma \'unica resolu\c{c}\~ao.} Tudo foi medido a $10{,}2$
    c\'elulas por raio. Um estudo de converg\^encia --- refinar $h$ e obter a
    taxa de cada m\'etrica --- n\~ao foi feito.
  \item \textbf{A vantagem do balan\c{c}o \'e espec\'ifica do caso est\'atico
    com $\kappa$ constante}, como a Se\c{c}\~ao~\ref{sec:elipse} mostrou: na
    elipse as duas formula\c{c}\~oes ficam equivalentes dentro da medida.
  \item \textbf{A curvatura na faceta \'e uma m\'edia ponderada dos marcadores
    vizinhos}, e n\~ao um campo de curvatura propriamente constru\'ido e
    estendido. Foi o suficiente para estabilizar a elipse, mas n\~ao foi testada
    em geometria com curvatura fortemente vari\'avel ou com mudan\c{c}a de
    sinal.
  \item \textbf{Serial.} $n_p=1$; o espalhamento paralelo n\~ao est\'a
    implementado.
  \item \textbf{Propriedades iguais.} A terceira mudan\c{c}a estrutural --- o
    salto de $\rho$ e $\mu$ --- est\'a adiada nos dois lados. Em escoamento com
    raz\~ao de densidade alta as conclus\~oes podem mudar.
  \item \textbf{A cirurgia foi exercitada} apenas na relaxa\c{c}\~ao da elipse,
    onde a frente encurta monotonicamente (108 para 102 marcadores). Esticamento
    forte, que exigiria \emph{inser\c{c}\~ao} sustentada, segue sem teste
    acoplado.
\end{itemize}

\subsection{Pr\'oximos passos}

Em ordem de valor:

\begin{enumerate}
  \item \textbf{Converg\^encia em $h$} das tr\^es m\'etricas, nos dois m\'etodos.
    D\'a \`a compara\c{c}\~ao o alcance que hoje lhe falta.
  \item \textbf{Fase B3}, oscila\c{c}\~ao de gota contra a frequ\^encia de Lamb:
    exercita din\^amica acoplada e, por deformar a interface, finalmente p\~oe a
    cirurgia \`a prova com o solver no circuito.
  \item \textbf{Paralelo}, com a redu\c{c}\~ao de franja no espalhamento.
\end{enumerate}

A compara\c{c}\~ao com o VOF pode acompanhar cada uma dessas etapas: o caso
existe, as sondas est\~ao instaladas nos dois lados, e as duas formula\c{c}\~oes
de for\c{c}a coexistem atr\'as de um seletor.


\bibliographystyle{plain}
\bibliography{referencias}

\end{document}
